\documentclass[xcolor={svgnames, table}, usepdftitle=false]{beamer}

% Beamer
\usetheme{Boadilla}
\usefonttheme{serif}
\setbeamertemplate{navigation symbols}{}
\setbeamertemplate{caption}[numbered]

% Fonts
\usepackage{fontspec}

% Languages
\usepackage{polyglossia}

% Bibliographies
\usepackage[hyperref, backend=biber]{biblatex}

% Microtype
\usepackage[babel]{microtype}

% Mathematics
\usepackage{mathtools}
\usepackage[math-style=TeX, bold-style=upright]{unicode-math}
\usepackage{algorithm}
\usepackage{algpseudocode}
\usepackage{prftree}

% Source code
\usepackage{fancyvrb}
\usepackage[newfloat]{minted}

% Logos
\usepackage{texlogos}

% Boxes
\usepackage{varwidth}
\usepackage[skins, minted]{tcolorbox}

% Tables
\usepackage{booktabs}
\usepackage{longtable}
\usepackage{tabu}

% Graphics and captioning
\usepackage{graphicx}
\usepackage[compatibility=false]{caption}
\usepackage{subcaption}

% Numbers
\usepackage[binary-units=true]{siunitx}
\usepackage{nth}

% Plotting
\usepackage{tikz}
\usepackage{pgfplots}
\usepackage{pgfplotstable}

% Cryptocode (load after `hyperref`)
\usepackage[primitives, operators, keys]{cryptocode}

% Quotations
\usepackage[strict]{csquotes}

% Dashes and hyphens
\usepackage[shortcuts]{extdash}

%
% Package setup
%

\setmainlanguage{english}
\setmathfont{Latin Modern Math}
\mathtoolsset{mathic=true}
\sisetup{%
  detect-all,
  detect-display-math,
  binary-units
}
\addbibresource{\jobname.bib}
\hypersetup{%
  unicode,
  pdfinfo={%
    Title={Trusted Timestamping},
    Subject={Computer Security},
    Author={Dario Gjorgjevski and Petar Tonkovikj},
    Keywords={Computer Security, Trusted Timestamping, Flask, OpenSSL}
  }
}
\usetikzlibrary{patterns, fit, positioning, calc, shapes.arrows}
\pgfplotsset{compat=newest}
\graphicspath{{Figures/}}

%
% Commands
%

% Vectors and matrices (boldface)
\renewcommand*{\vec}{\symbf}
\newcommand*{\mat}{\symbf}

% Expected value and variance
\DeclareMathOperator{\E}{\mathbb{E}}
\DeclareMathOperator{\var}{var}

% Matrix rank
\DeclareMathOperator{\rank}{rank}

%
% Title
%

\title{Notary}
\subtitle{A Flask Application for Trusted Timestamping}
\author[Dario Gj.\ \& Petar T.]{%
  Dario Gjorgjevski\inst{1} \and Petar Tonkovikj\inst{1}
}
\institute[FCSE]{%
  \inst{1}Faculty of Computer Science and Engineering\\%
  Ss.\ Cyril and Methodius University in Skopje
}
\date{\today}

%
% Logo
%

\logo{\includegraphics[height=0.66cm]{../Logo_ФИНКИ.pdf}}

%
% Table of contents
%

\AtBeginSection[]{%
  \begin{frame}{Contents}
    \tableofcontents[currentsection, hideothersubsections]
  \end{frame}
}

%
% Document
%

\begin{document}

\begin{frame}
  \titlepage{}
\end{frame}

\section{Introduction}

\begin{frame}{Trusted Timestamping}
  \begin{itemize}
  \item Trusted timestamping is the process of securely keeping track
    of the creation time of a document.
  \item \enquote{Securely} roughly means that \emph{no one}\--- not
    even the owner of the document\--- should be able to change it.
  \item \href{https://tools.ietf.org/html/rfc3161}{RFC3161} specifies
    a scheme for trusted timestamping for the PKI \(\implies\) OpenSSL
    provides an open\-/source implementation of RFC3161.
  \end{itemize}
\end{frame}

\begin{frame}{Web Application}{Architectural choices}
  \begin{itemize}
  \item Must ensure independent timestamp generation and verification
    \(\implies\) two parts meant to be run in isolated environments.
  \item Both parts are built using
    Flask\footnote{\url{http://flask.pocoo.org/}} and Python 3.
  \item Simple and bulletproof architecture built on top of the Web
    Server Gateway Interface (WSGI).
  \end{itemize}
\end{frame}

\begin{frame}{Web Application}{Deployment}
  \begin{itemize}
  \item WSGI scripts can be deployed on the Apache HTTP Server via
    \Verb+mod\_wsgi+.
  \item HTTPS and user authentication are provided by OpenSSL via
    \Verb+mod\_ssl+.
  \item The certificate used for SSL is signed by our TA, Hristina.
  \item The certificate used for trusted timestamping is generated and
    self\-/signed by us.
  \item We allow only users with certificates issued by our
    organization to authenticate themselves to the website.  See
    listing~\ref{lst:authentication}.
  \end{itemize}
\end{frame}

\begin{frame}[fragile]{Web Application}{Authentication}
\begin{listing}[H]
\centering
\begin{minted}[fontsize=\small]{aconf}
#   Access Control:
#   With SSLRequire you can do per-directory access control based
#   on arbitrary complex boolean expressions containing server
#   variable checks and other lookup directives.  The syntax is a
#   mixture between C and Perl.  See the mod_ssl documentation
#   for more details.
<Location>
SSLRequire %{SSL_CLIENT_I_DN_O} in \
           {"Faculty of Computer Science and Engineering", \
	    "Ss. Cyril and Methodius University", \
	    "FCSE", "UKIM"}
</Location>
\end{minted}
\caption{OpenSSL Access Control}\label{lst:authentication}
\end{listing}
\end{frame}

\section{Algorithm for Trusted Timestamping}

\begin{frame}{Generating a Trusted Timestamp}{Informally}
  Generating a trusted timestamp is straightforward and consists of:
  \begin{enumerate}
  \item Hashing the data that needs to be timestamped;
  \item Concatenating the hash with the timestamp;
  \item Hashing the concatenation of the two;
  \item Digitally signing the hash with the private key of the TSA.\@
  \end{enumerate}
\end{frame}

\begin{frame}{Generating a Trusted Timestamp}{Formally}
  A trusted timestamp contains \texttt{data}, a \texttt{timestamp}, a
  hash \texttt{H}, and a signature \texttt{S}.  The latter two are defined as:
  \begin{equation}
    \begin{aligned}
      \mathtt{H} &\coloneqq \hash(\hash(\mathtt{data}) \concat \mathtt{timestamp})\text{, and} \\
      \mathtt{S} &\coloneqq \sig_{\mathtt{TSA}}(\mathtt{H})\text,
    \end{aligned}
  \end{equation}
  where \(\sig_{\pk}(\cdot)\) is a digital signature under the private
  key \(\pk\), \(\hash(\cdot)\) is a cryptographic hash function such
  as SHA\=/\num{256} or SHA\=/\num{512}, and \enquote{\(\concat\)}
  denotes concatenation.
\end{frame}

\begin{frame}{Verifying a Trusted Timestamp}
  \begin{enumerate}
  \item Hash data and tentative timestamp as previously.
  \item Verify the TSA's signature on the provided hash.
  \item Compare the computed hash and the provided hash.
  \end{enumerate}

  The two hashes being identical is a strong guarantee of the
  timestamp being correct (i.e., issued by the TSA) and unaltered.
\end{frame}

\section{Workflow}

\begin{frame}{Supplying the Authentication Certificate}
  All major web browsers are able to supply PKI certificates for
  authentication to the website.  An example doing this on our website
  is found on figure~\ref{fig:certificate-prompt}.
  \begin{figure}[H]
    \centering
    \includegraphics[width=.5\textwidth]{CertificatePrompt.png}
    \caption{Mozilla Firefox Prompting for Certificate}\label{fig:certificate-prompt}
  \end{figure}
\end{frame}

\begin{frame}{Verifying the SSL Connection}
  \begin{figure}[H]
    \centering
    \includegraphics[width=.75\textwidth]{SecureConnection.png}
    \caption{SSL Certificate Information}\label{fig:secure-connection}
  \end{figure}
\end{frame}

\begin{frame}{Uploading Documents for Trusted Timestamping}
  \begin{itemize}
  \item Uploading is as simple as pressing the \beamerbutton{Upload}
    button.
  \item Once on the server, document can be either:
    \begin{itemize}
    \item Timestamped; or
    \item Deleted.
    \end{itemize}
  \item Pressing the \beamerbutton{Timestamp} button sends to the user
    a zip file with a:
    \begin{itemize}
    \item Timestamp query (TSQ) used to timestamp the file; and a
    \item Timestamp response (TSR) containing the trusted timestamp.
    \end{itemize}
  \end{itemize}
\end{frame}

\begin{frame}{Uploading the TSQ and TSR for Verification}
  The TSQ and TSR files downloaded from the file server can be
  uploaded to the second web application for verification.  Once
  uploaded, they are passed as input to the \mintinline{bash}{ts
    -verify} OpenSSL command.
  \begin{figure}[H]
    \includegraphics[width=.75\textwidth]{UploadTrustedTimestamp.png}
    \caption{Uploading the Required Files}
  \end{figure}
\end{frame}

\begin{frame}{Inspecting the Verification Results}
  Once a user has uploaded the TSQ and TSR files, he/she will be
  redirected to a page containing the results of the verification.
  \begin{figure}[H]
    \includegraphics[width=.75\textwidth]{VerificationResults.png}
    \caption{Output of \mintinline{bash}{ts -verify}}
  \end{figure}
\end{frame}

\begin{frame}{Pros and Cons of Having Two Separate Applications}
  \begin{columns}
    \begin{column}{.55\textwidth}
      \tcbox[title=Pros, colframe=DarkGreen!75!white]{%
        \begin{minipage}[t][.6\textheight][t]{\dimexpr.75\textwidth-2\fboxsep-2\fboxrule-5pt\relax}
          \setlength{\leftmargini}{2pt}
          \begin{itemize}
          \item Users to not need to use OpenSSL at all \(\implies\) more
            user\-/friendliness and less error\-/proneness.
          \item Separate entities for timestamping and verification
            \(\implies\) more trustworthiness since the compromise of one
            entity can be \enquote{detected} by the other.
          \end{itemize}
        \end{minipage}
      }
    \end{column}
    \begin{column}{.4\textwidth}
        \tcbox[title=Cons, colframe=Red!75!white]{%
          \begin{minipage}[t][.6\textheight][t]{\dimexpr.75\textwidth-2\fboxsep-2\fboxrule-5pt\relax}
            \setlength{\leftmargini}{2pt}
            \begin{itemize}
            \item Deployment: it is harder to deploy and keep two
              applications up\=/to\=/date.
            \item Convenience: users need to use two different
              applications for a single functionality.
            \end{itemize}
          \end{minipage}
        }
    \end{column}
  \end{columns}
\end{frame}

\section{Conclusion and Future Work}

\begin{frame}{Conclusion}
  \begin{itemize}
  \item We developed a simple trusted timestamping server.
  \item Suitable for deployment in \emph{closed} environments: closed
    in a sense that users can authenticate themselves via
    certificates.
  \item All components utilized are open\-/source, simple, and
    have been used and studied extensively.
  \end{itemize}
\end{frame}

\begin{frame}{Future Work}
  \begin{itemize}
  \item Store files on a database for easier fault\-/tolerance and
    backups.  At present, files are stored on the file system as\=/is.
  \item Functionality to timestamp files the moment they are uploaded
    rather than upon request.
  \item Alternative ways of authentications (user accounts).
  \item Private files.
  \item Same certificate for the TSA and the SSL connection.
  \end{itemize}
\end{frame}

\end{document}

%%% Local Variables:
%%% mode: latex
%%% TeX-command-extra-options: "-shell-escape"
%%% TeX-engine: luatex
%%% TeX-master: t
%%% End: